\section[On the homotopy type of the space of almost complex structures on six-manifolds]{On the homotopy type of the space\\ of almost complex structures on six-manifolds}

In this section we study the homotopy type of the components of the space of almost complex structures on connected six-manifolds $M$, complementing our previous treatment~\cite{FGM21} of the case when $c_1=0$.

\begin{Theorem}\label{fibration}
Let $M$ be an oriented Riemannian manifold of dimension four or six, $J$ an orthogonal almost complex structure on $M$ compatible with the orientation, and $J(M)$ its component in the space of orthogonal almost complex structures. Let $S_{\C}^+(M)$ be the positive spinor bundle on $M$ determined by $J$. Then the map
\[
\Gamma\big( S\big(S_\C^+(M)\big)\big) \to J(M)
\]
is surjective, and for any lift $s$ of $J$,
\[
\Map\big(M,S^1\big) \to \Gamma\big( S\big(S_\C^+\big)\big)_s \to J(M)
\]
is a fiber sequence.
\end{Theorem}
\begin{proof}
Identifying $S_\C^+(M)$ with $\oplus_{k \text{ even}} \Lambda^k_\C(TM)$, Remark~\ref{lift} shows that
the constant section $s \colon M \to \Lambda^0_\C(TM) \subset S_\C^+(M)$ defined by
$s(x)=1$ lifts the section $J$. As $S\big(S_\C^+\big) \to Z_+(M)$ is a~fiberwise fibration over $M$, the map induced on sections is a~fibration. Since the fiber over $J$ is nonempty, any choice of lift of $J$ identifies the
fiber with $\Map\big(M,S^1\big)$.
\end{proof}

\begin{Remark} The above argument remains valid in all dimensions if one replaces $S\big(S_{\C}^+(M)\big)$ with $S\big(S_{\C}^+(M)\big) \cap {\mathcal P}S ^+_{\C}(M)$. \end{Remark}

Now let $(M,J)$ be a closed almost complex Riemannian six-manifold with $H^1(M;\ZZ) = 0$ (equivalently, $b_1 = 0$) satisfying $c_1c_2 - c_3 \neq 0$. By Theorem~\ref{fibration}, we have the fibration $\Map\big(M,S^1\big) \to \Gamma\big( S\big(S_\C^+(M)\big)\big)_s \to J(M)$. Since $H^1(M;\ZZ) = 0$, evaluation at a chosen basepoint gives a homotopy equivalence $\Map\big(M,S^1\big) \xrightarrow{{\rm ev}} S^1$, so we have a principal fiber sequence $S^1 \to \Gamma\big(S\big(S_\C^+(M)\big)\big)_s \to J(M)$.

We can thus consider instead the fiber sequence \begin{equation}\label{fibration2} \Gamma(S\big(S_{\C}^+\big))_s \to J(M) \to {\mathcal B}S ^1 \simeq \CP^\infty. \end{equation} Now, $S\big(S_{\C}^+\big)$ is an oriented fiber bundle over $M$ with fiber $S^7$, and hence it is classified by a map $M \to {\mathcal B}{\rm Aut}^+\big(S^7\big)$ to the classifying space of orientation-preserving homotopy automorphisms of~$S^7$. It is known that ${\mathcal B}{\rm Aut}^+\big(S^7\big)_{\QQ} \simeq K(\QQ, 8)$, where the subscript of $\QQ$ denotes (Sullivan) rationalization~\cite[Section~11]{S77}. By~\cite[Theorem~5.3]{M87}, the rationalization of $\Gamma\big(S\big(S_{\C}^+\big)\big)_s$ is homotopy equivalent to (the connected component corresponding to $s$ in) the space of sections of the fiberwise rationalized $S^7$ bundle over $M$. Since $H^8(M;\QQ) = 0$, this latter bundle is trivial, and hence $\big(\Gamma\big(S\big(S_{\C}^+\big)\big)_s\big)_{\QQ}$ is homotopy equivalent to a connected component of $\Map\big(M, S^7_{\QQ}\big)$. Denoting the Betti numbers of $M$ by $b_i$, by Thom's theorem on the space of maps into an Eilenberg--Maclane space, $\Map\big(M, S^7_{\QQ}\big)$ is in fact connected and has the homotopy type of $S^1_{\QQ} \times S^7_{\QQ} \times K(\QQ, 3)^{b_4} \times K(\QQ, 4)^{b_3} \times K(\QQ,5)^{b_2}$.

We can describe the fundamental group of $J(M)$ by using a theorem of Crabb and Sutherland:

\begin{Theorem}[{\cite[Theorem~2.12(i)]{CS84}}] Let $M$ be a closed connected $2n$-manifold and $\xi$ a complex $(n+1)$-plane bundle over $M$. Denote by $N\xi$ any component of the space of sections of $\mathbb{P}(\xi)$ whose elements lift to sections of $\xi$. Then $\pi_1(N\xi)$ is a central extension \[0 \to \ZZ/\big(\big\langle c_n(\xi), [M] \big\rangle \big) \to \pi_1(N\xi) \to H^1(M) \to 0.\]
\end{Theorem}

By Theorem~\ref{fibration}, we can apply this to $\xi = S_{\C}^+(M) \cong \C \oplus \Lambda_{\C}^2 TM$, which satisfies ${N \xi = J(M)}$. Therefore, the fundamental group of $J(M)$ is given by $\ZZ$ modulo $\int_M\! c_1c_2-c_3$ (where ${c_i \! =\! c_i(TM)}$), see Corollary~\ref{cohomtwis}. By~\cite{M87}, $J(M)$ is a nilpotent space as it is the space of sections of a fibration with nilpotent fiber \big(namely $\CP^3$\big) over a finite-dimensional base. We may thus rationalize the fibration (\ref{fibration2}) to obtain the fibration \[\big(\Gamma\big(S\big(S_{\C}^+\big)\big)_s\big)_{\Q} \to J(M)_{\Q} \to K(\Q, 2).\] Consider the degree one rational class corresponding to the factor $S^1_{\QQ}$ in the fiber. If $c_1c_2 - c_3 \neq 0$, then since $\pi_1(J(M)_\Q) = 0$, this class must hit (a nonzero multiple of) the degree two generator in $K(\Q,2)$ in the Serre spectral sequence. A model for the total space is given by the tensor product of the base and fiber's models with a perturbed differential on the fiber generators~\cite[Section~4]{S77}; the previous sentence thus tells us that a model for $J(M)$ is of the form
\[\big(\Lambda\big(x_2, z_1, z_7, z_3^{i}, z_4^{j}, z_5^{k}\big),\, {\rm d}z_1 = x_2, \, {\rm d}z_7 \in (x_2), \, {\rm d}z_3^i \in (x_2), \, {\rm d}z_4^j \in (x_2),\, {\rm d}z_5^k \in (x_2)\big),
\] where $(x_2)$ denotes the ideal generated by $x_2$, and $i$, $j$, $k$ range over sets of size $b_4$, $b_3$, $b_2$, respectively.

Notice that such a model is not minimal, and in fact a minimal model is obtained by quotienting out the differential ideal generated by~$z_1$. Indeed, by the argument in~\cite[Proposition~2]{VS76} we see that the map $(\Lambda, {\rm d}) \to \big(\Lambda/(z_1, {\rm d}z_1)\Lambda, \bar{{\rm d}}\big)$ induces an isomorphism on cohomology.
Here by $(\Lambda,{\rm d})$ we denote the differential graded algebra displayed above, and by $\bar{{\rm d}}$ the induced differential on the quotient. We see that
\[
(\Lambda/(z_1, x_2)\Lambda, \bar{{\rm d}}) \cong \big(\Lambda\big(z_7, z_3^{i}, z_4^{j}, z_5^{k}\big), \bar{{\rm d}}=0\big),
\] where the right-hand side is minimal. To summarize, we have the following:

\begin{Theorem}\label{Qhomotopy6} Let $M$ be a connected closed six-manifold with $b_1 = 0$ equipped with an almost complex structure $J$ such that $\int_M c_1c_2-c_3 \neq 0$. Then the space of almost complex structures on~$M$ in the component of $J$ is a nilpotent space with finite cyclic fundamental group and minimal model given by $\big(\Lambda\big(z_7, z_3^{i}, z_4^{j}, z_5^{k}\big), {\rm d} = 0\big)$; here $i$, $j$, $k$ range over sets of size $b_4(M)$, $b_3(M)$, $b_2(M)$ $(=b_4(M))$, respectively. In particular, this space is formal.
\end{Theorem}

We emphasize that the above makes no nilpotency assumption on $M$. Alternatively, a model for (a given component of) the space of almost complex structures can be obtained with the Haefliger--Sullivan model for the space of sections of a fibration~\cite[Section~11]{S77},~\cite{H82}, applied to $\CP^3 \to Z_+(M) \to M$. Our approach above lets one quickly describe the minimal model of~$J(M)$ using the geometry of the setup.
